\chapter{Materiais e Métodos}
\label{cap:metodologia}

Este capítulo descreve o ambiente físico, as ferramentas de software,
a abordagem metodológica e a divisão de responsabilidades adotadas
durante o desenvolvimento do infográfico.

\section{Hardware utilizado no desenvolvimento}
\label{sec:hardware}

O artefato foi desenvolvido em um \textbf{notebook Acer Aspire A315-42}
com configuração de uso típico de aluno de graduação, conforme
detalhado na Tabela~\ref{tab:hw}.

\begin{table}[h!]
\centering
\caption{Especificações do equipamento de desenvolvimento}
\label{tab:hw}
\begin{tabular}{lp{8.5cm}}
\toprule
\textbf{Componente} & \textbf{Especificação} \\
\midrule
Processador & AMD Ryzen 5 3500U, 4 \emph{cores} físicos / 8
              \emph{threads}, 2,1\,\textsc{ghz} base \\
Arquitetura  & x86\_64, AMD Zen+, fabricação 12\,\textsc{nm} \\
Memória RAM  & 9,6\,\textsc{gib} \textsc{ddr4} 2400\,\textsc{mhz}
               (não-\textsc{ecc}, conforme classe consumidor) \\
Armazenamento & SSD \textsc{nvme} M.2, 512\,\textsc{gb} \\
Sistema operacional & EndeavourOS (base Arch Linux),
                      \emph{kernel} 7.0.9-zen1 \\
Navegador alvo & Chromium 134+, Mozilla Firefox 130+ \\
\bottomrule
\end{tabular}
\end{table}

Por se tratar de aplicação web client-side, o infográfico foi
testado também em equipamentos diversos dos demais integrantes da
equipe, com sistemas operacionais Windows~11 e macOS~14, confirmando
compatibilidade transversal nos navegadores baseados em Chromium e
no Mozilla Firefox.

\section{Ferramentas e ambiente de desenvolvimento}
\label{sec:ferramentas}

\subsection{Editor de código}

O desenvolvimento foi conduzido no \textbf{Visual Studio Code}
\cite{vscode2026}, distribuído pela Microsoft sob licença
\textsc{mit}. As extensões utilizadas restringiram-se ao conjunto
mínimo necessário: \textit{Prettier} para formatação consistente,
\textit{Live Server} para auto-reload durante o desenvolvimento e
\textit{GitLens} para visualização do histórico de commits.

\subsection{Assistentes de IA generativa}

A produção utilizou duas ferramentas de IA, declaradas conforme a
boa prática acadêmica atual de transparência sobre uso de modelos
generativos:

\begin{description}
    \item[Google Gemini --- modo Canvas~\cite{gemini2026}.] Empregado
          na fase inicial de prototipagem, especialmente na geração
          de andaime visual em HTML/CSS para os módulos
          interativos (estrutura inicial do simulador \textsc{raid} e
          das barras de recursos da virtualização).
    \item[Anthropic Claude Code~\cite{claudecode2026}.] Empregado na
          segunda iteração para aprofundamento técnico do conteúdo,
          refatoração da disposição em sub-abas, escrita das tabelas
          comparativas, redação deste relatório e elaboração do
          roteiro de defesa. A interação se deu via terminal sobre
          o diretório do projeto, com revisão humana de cada commit.
\end{description}

Ambas as ferramentas atuaram como copiloto: o discernimento técnico,
a verificação dos números citados e a aderência ao material da
disciplina foram conduzidos pelos integrantes da equipe.

\subsection{Versionamento e hospedagem}

O código-fonte está versionado em \textsf{Git} e hospedado no
GitHub no repositório
\url{https://github.com/glkwolff/Projeto-Estruturas-de-computadores}.

\subsection{Stack tecnológico do artefato}

O infográfico final é \emph{auto-contido}: um único arquivo
\texttt{projeto.html} (\textasciitilde 2.500 linhas) contendo
estrutura, estilos e lógica. Não há \emph{build step}, transpilação
ou dependência de pacotes externos. A Tabela~\ref{tab:stack} resume.

\begin{table}[h!]
\centering
\caption{Tecnologias usadas no artefato final}
\label{tab:stack}
\begin{tabular}{lp{8cm}}
\toprule
\textbf{Tecnologia} & \textbf{Uso} \\
\midrule
HTML5            & Estrutura semântica em \texttt{section}/
                   \texttt{article}, \texttt{table} \\
CSS3             & \emph{Custom properties}, \texttt{grid},
                   \texttt{flex}, \texttt{clip-path}, transições,
                   \emph{media queries} \\
JavaScript ES6+  & Manipulação de DOM, lógica de estado em
                   variáveis-módulo, \texttt{setTimeout} para
                   animações, \texttt{const}/\texttt{let} \\
SVG inline       & (Não utilizado --- toda visualização é
                   feita por divs estilizadas, para reduzir
                   dependência) \\
\bottomrule
\end{tabular}
\end{table}

A escolha de \emph{vanilla} JavaScript --- sem frameworks como React
ou Vue --- foi deliberada: o objetivo pedagógico é exibir conceitos
de arquitetura, não exercitar engenharia de \emph{front-end}.
Manter o projeto em arquivo único também simplifica entrega,
avaliação e execução: basta abrir o HTML no navegador.

\section{Metodologia de desenvolvimento}
\label{sec:metodologia-dev}

O projeto seguiu duas iterações sucessivas:

\begin{enumerate}
    \item \textbf{Iteração 1 (abril/2026)}. Versão inicial entregue
          como parcial, focada em estabelecer o esqueleto visual e os
          cinco módulos básicos com analogias e curiosidades. Avaliada
          como insuficientemente profunda pelo orientador.
    \item \textbf{Iteração 2 (junho/2026)}. Refatoração com foco em
          \emph{profundidade conteudística}: incorporação de tabelas
          quantitativas, fórmulas explícitas, simuladores adicionais
          (calculadora \textsc{raid}, simulador de cache hit/miss,
          escalonador rubro-negro), reorganização de cada módulo em
          sub-abas para reduzir densidade visual e adição de um
          sexto módulo (\textit{stack} completo) integrador.
\end{enumerate}

A metodologia em cada módulo seguiu o ciclo:

\begin{enumerate}
    \item leitura do material das aulas correspondente
          (aulas 6 a 12);
    \item mapeamento dos conceitos para subseções navegáveis;
    \item definição de uma interação central ---
          ``o que o aluno pode \emph{fazer} aqui?'' ---
          em vez de ler texto passivo;
    \item implementação progressiva, testando hot-reload no
          navegador;
    \item revisão pela equipe e ajuste fino.
\end{enumerate}

\section{Divisão de responsabilidades na equipe}
\label{sec:equipe}

A equipe é composta por cinco integrantes, com divisão temática
para garantir profundidade individual em pelo menos um módulo
(condição da nota de defesa, que é individual). A
Tabela~\ref{tab:equipe} resume.

\begin{table}[h!]
\centering
\caption{Distribuição de responsabilidades por integrante}
\label{tab:equipe}
\begin{tabular}{ll}
\toprule
\textbf{Integrante} & \textbf{Módulo principal de apresentação} \\
\midrule
Gabriel Wolff       & Módulo 1 --- Desktop \textit{vs.}\ Servidor \\
Vitor Camillo       & Módulo 2 --- RAID + Lei de Amdahl \\
João Victor Alvarez & Módulo 3 --- Memória \textsc{ecc} + Hamming \\
Lucas Ashikaga      & Módulo 4 --- Hierarquia \& \textsc{numa} \\
Wellington Fonseca  & Módulo 5 --- Virtualização \\
Equipe (coletivo)   & Módulo 6 --- \emph{Stack} completo (síntese) \\
\bottomrule
\end{tabular}
\end{table}

A divisão temática não impediu colaboração em todos os módulos: cada
integrante revisou e contribuiu para os demais, mas tem domínio
aprofundado do tema que apresentará na defesa.

\section{Critérios de avaliação e entrega}

Conforme orientação do professor, a avaliação é composta de duas
notas independentes:

\begin{itemize}
    \item \textbf{Apresentação de defesa} (individual, 2,5 pontos):
          duração máxima de 15 minutos, com cada integrante
          apresentando seu módulo principal.
    \item \textbf{Relatório escrito} (coletivo, 2,5 pontos): este
          documento, com estrutura \textsc{abnt}.
\end{itemize}

A entrega final compreende um arquivo \texttt{.zip} contendo o
relatório em PDF, a apresentação, o arquivo \texttt{projeto.html} do
infográfico e o \emph{link} para o repositório GitHub (uma vez que
o código-fonte completo ultrapassa as 2\,500 linhas, conforme
preconizado pelo critério de entrega).
